\section{Part~1 core: exactness and isotypic detectability (minimal statement)}
\label{sec:part1core}

\paragraph{Status of this section.}
The full Part~1 (workstreams N1-1 through N1-4: unfolded half-period shift,
cross-cycle exchangeability, sequential e-processes, IPM refinement of the
$\varepsilon$-budget) is not yet in this document. Part~2
(Section~\ref{sec:part2}) needs exactly two of its results in a minimal form
--- the exactness of the flip test under $H_0$ and the isotypic
characterization of what an invariance test can see --- and states them here
with proofs, so that Part~2 is self-contained. When the full Part~1 lands,
Theorems~\ref{thm:n1-1} and~\ref{thm:n1-2} are to be replaced by their general
versions and the labels kept.
\todo[inline]{Part~1 stand-in (Sprint~6): replace by the full N1-1/N1-2
statements (unfolded shift, e-processes, IPM route) and re-point the labels
\texttt{thm:n1-1}, \texttt{thm:n1-2}.}

\subsection{Setting}
\label{sec:p1-setting}

Cycle elements live in the finite-dimensional space $\mathcal{W} = \R^{dN}$
($d$ channels on the $N$-point phase grid, footnote of
Section~\ref{sec:cycles}). For the trot, $\symgrp = \{(e,0),\sigma_*\}$ with
$\sigma_* = (\gs,\tfrac12)$, and $\rho := \rho(\sigma_*)$ acts on
$\mathcal{W}$ by \eqref{eq:cycle-action}: a signed channel permutation
composed with the half-grid cyclic shift. Both factors are orthogonal
involutions, hence
\begin{equation}
  \label{eq:rho-involution}
  \rho^{2} = I, \qquad \rho^{\top} = \rho, \qquad
  \Pi^{\pm} := \tfrac12 (I \pm \rho), \qquad
  \mathcal{W} = \mathcal{W}^{+} \oplus \mathcal{W}^{-},
\end{equation}
where $\Pi^{\pm}$ are complementary orthogonal projectors ---
the isotypic decomposition of $\mathcal{W}$ into the trivial ($\rho = +1$)
and sign ($\rho = -1$) components of $\symgrp \cong \mathbb{Z}_2$.
Throughout this section the cycles are assumed independent (or, more weakly,
exchangeable) across $k$ --- the separate requirement of
Remark~\ref{rem:exch}; how far the wrapped element satisfies the invariance
exactly is the matter of Remark~\ref{rem:wrap} and is not re-litigated here.

\subsection{Theorem N1-1: exactness of the flip test}
\label{sec:n1-1}

Let $T : \mathcal{W}^{K} \to \R$ be any statistic of $K$ cycle elements. The
flip group $\mathcal{G}_K := \{ \rho^{\varepsilon_1} \times \dots \times
\rho^{\varepsilon_K} : \varepsilon \in \{0,1\}^{K}\}$ acts on
$(Z_1,\dots,Z_K)$ coordinatewise. The random-subset test of
\cite{hemerik2018exact} draws $g_2,\dots,g_M$ i.i.d.\ uniformly from
$\mathcal{G}_K$, sets $g_1 = \mathrm{id}$, and reports
\begin{equation}
  \label{eq:hg-p}
  p \;=\; \frac{1 + \#\{ j \ge 2 : T(g_j Z) \ge T(Z) \}}{M}.
\end{equation}

\begin{theorem}[N1-1, exactness]
\label{thm:n1-1}
If $Z_1,\dots,Z_K$ are independent and each satisfies $H_0$, that is
$\Law(\rho Z_k) = \Law(Z_k)$, then $\Law(g Z) = \Law(Z)$ for every
$g \in \mathcal{G}_K$, and the $p$-value \eqref{eq:hg-p} is super-uniform:
$\mathbb{P}(p \le \alpha) \le \alpha$ for all $\alpha \in [0,1]$, with
equality at $\alpha \in \{1/M, 2/M, \dots\}$ when $T$ has no ties. The
statement holds for every statistic $T$ and does not involve any model of the
data.
\end{theorem}

\begin{proof}
Independence and per-cycle invariance give
$\Law(\rho^{\varepsilon_1} Z_1, \dots, \rho^{\varepsilon_K} Z_K) =
\bigotimes_k \Law(\rho^{\varepsilon_k} Z_k) = \bigotimes_k \Law(Z_k)$, i.e.\
$\Law(gZ) = \Law(Z)$ for all $g \in \mathcal{G}_K$. Given this invariance,
the multiset $\{T(g_1 Z), \dots, T(g_M Z)\}$ with $g_1 = \mathrm{id}$ and
$g_2,\dots,g_M$ uniform is exchangeable: for the joint law of
$(g_1 Z, g_2 Z, \dots, g_M Z)$, replacing $Z$ by $g Z$ ($g$ uniform,
independent) leaves the law of $Z$ unchanged and maps the tuple to
$(g Z, g_2 g Z, \dots)$ whose entries $g_j g$ are again i.i.d.\ uniform, so
the identity holds no special position. The rank of $T(Z)$ among $M$
exchangeable values is uniform on $\{1,\dots,M\}$ (ties resolved
conservatively), which is
\cite[Theorem~2]{hemerik2018exact}; \eqref{eq:hg-p} is that rank divided by
$M$.
\end{proof}

\begin{remark}[What the theorem does not say]
It says nothing about power, and it presumes cross-cycle independence
(exchangeability suffices, block flips handle short-range dependence,
Remark~\ref{rem:exch}). Under the $\varepsilon$-approximate assumptions the
level is inflated additively by Lemma~\ref{lem:eps}. It is also the reason
why every downstream statistic in the project --- including the residual
statistics of Part~2 --- is built from $\rho$-paired quantities: the theorem
holds for any $T$, so the choice of $T$ is a pure power question.
\end{remark}

\subsection{Theorem N1-2: isotypic detectability}
\label{sec:n1-2}

Structured alternatives (Definition~\ref{def:h1}) shift the mean trajectory.
The minimal statement is for location alternatives with $\symgrp$-invariant
fluctuations; it is what Part~2 transports to residuals.

\begin{theorem}[N1-2, isotypic detectability]
\label{thm:n1-2}
Let $Z_k = \mu + \xi_k$ with $\mu \in \mathcal{W}$ deterministic and
$\xi_k$ i.i.d.\ with $\Law(\rho \xi_k) = \Law(\xi_k)$ and
$\mathbb{E}\xi_k = 0$, $\operatorname{Var}(\xi_k) = \Sigma$. Then:
\begin{enumerate}
  \item[(i)] \emph{(invisible component)} If $\Pi^{-}\mu = 0$ then $H_0$
  holds for $(Z_k)$, and \emph{no} $\symgrp$-invariance test --- no test
  whose level is controlled by Theorem~\ref{thm:n1-1} --- has power above its
  level against $\mu$: the trivial component $\Pi^{+}\mu$ is structurally
  invisible.
  \item[(ii)] \emph{(visible component)} If $\Pi^{-}\mu \neq 0$ then $H_0$
  fails, and the paired-difference statistic
  $T(Z) = \|\tfrac1K \sum_k D_k\|^2$, $D_k := Z_k - \rho Z_k = 2\Pi^{-}Z_k$,
  has $\mathbb{E}\,T = 4\|\Pi^{-}\mu\|^2 + \tfrac{1}{K}\operatorname{tr}
  \operatorname{Var}(D_1)$ while its flip-null counterpart (any nontrivial
  flip pattern) has expectation $\tfrac{1}{K}\operatorname{tr}\operatorname{Var}(D_1)
  + O(\|\Pi^-\mu\|^2/K)$: the flip test is consistent as $K \to \infty$, with
  noncentrality growing like $K\,\|\Pi^{-}\mu\|^{2}$ (in the standardized
  version, $K \sum_c (\Pi^-\mu)_c^2 / \operatorname{Var}(D_1)_{cc}$).
  \item[(iii)] \emph{(general alternatives)} An arbitrary alternative law
  $P_1$ on $\mathcal{W}$ is undetectable by every $\symgrp$-invariance test
  iff $\push{\rho} P_1 = P_1$; in particular a fault pattern fixed by
  $\sigma_*$ (bilateral, synchronized, equal magnitude,
  Example~\ref{ex:bilateral}) is invisible to $R^{-}$ whatever its magnitude,
  and must be read by a magnitude ($\Pi^{+}$) channel.
\end{enumerate}
\end{theorem}

\begin{proof}
(i) $\Pi^{-}\mu = 0$ means $\rho\mu = \mu$, so $\rho Z_k = \mu + \rho\xi_k
\eqdist \mu + \xi_k = Z_k$: $H_0$ holds and Theorem~\ref{thm:n1-1} makes every
flip test level-$\alpha$ under $\Law(Z)$, i.e.\ its rejection probability is
at most $\alpha$. (ii) $D_k = 2\Pi^{-}\mu + (\xi_k - \rho\xi_k)$ with
$\mathbb{E}(\xi_k - \rho\xi_k) = 0$; hence $\mathbb{E}\|\bar D\|^2 =
\|2\Pi^-\mu\|^2 + \operatorname{tr}\operatorname{Var}(D_1)/K$. For a flip
pattern $\varepsilon$ with $m$ of the $K$ signs equal to $-1$, the flipped
mean is $\tfrac{1}{K}\sum_k \varepsilon_k D_k$ whose expectation is
$\tfrac{K-2m}{K}\,2\Pi^-\mu$, of squared norm $\le 4\|\Pi^-\mu\|^2$ with
equality only for the identity; averaged over uniform flips the mean of
$\tfrac{K-2m}{K}$ is $0$ and its second moment is $1/K$, giving the stated
$O(\|\Pi^-\mu\|^2/K)$ term. Consistency follows because the observed
statistic concentrates at $4\|\Pi^-\mu\|^2 > 0$ while all but a vanishing
fraction of flipped statistics concentrate at $0$ (law of large numbers in
both cases). If $\Pi^-\mu \ne 0$ then $\rho\mu \ne \mu$ and, since
$\mathbb{E}\rho Z = \rho\mu \neq \mu = \mathbb{E}Z$, $H_0$ fails.
(iii) If $\push{\rho}P_1 = P_1$ the data satisfy $H_0$ and (i) applies; if
$\push{\rho}P_1 \neq P_1$ some bounded measurable $h$ has
$\mathbb{E}_{P_1}[h(Z) - h(\rho Z)] \neq 0$, and the flip test with
$T = |\tfrac1K \sum_k (h(Z_k) - h(\rho Z_k))|$ is consistent by the argument
of (ii).
\end{proof}

\begin{remark}[Reading]
Detection lives in the sign component $\mathcal{W}^-$, blindness in the
trivial component $\mathcal{W}^+$; the two are orthogonal complements. This
is the structural reason for the two-channel design ($R^-$ invariance test
plus a magnitude channel), which Part~2 sharpens: with an equivariant nominal
model both channels read the \emph{same} residual, one component each.
The empirical anchor is e04c (rp005): the bilateral-equal group is $R^-$-blind
with window rejection in the binomial band, the single-leg group is detected
by both channels, and the unequal group is intermediate in the antisymmetric
\emph{share}; the observed superlinearity of the response in $1-\kappa$
(rp005, Finding) is a reminder that (ii) is a first-order statement in the
mean shift.
\end{remark}
